\documentclass[11pt,a4paper]{article}
\usepackage[utf8]{inputenc}
\usepackage[T1]{fontenc}
\usepackage[french]{babel}
\usepackage{lmodern}
\usepackage{amsmath,amssymb}
\usepackage[margin=2.5cm]{geometry}
\usepackage{microtype}
\usepackage{hyperref}
\usepackage{listings}
\lstset{basicstyle=\ttfamily\small, frame=single, breaklines=true,
  columns=fullflexible, keepspaces=true}
\title{Projet MLSF : résolution du jeu Meowdoku\\
  \large Question 2 : lecture DIMACS et algorithme DPLL}
\author{}
\date{}

\begin{document}
\maketitle

\section{Objectif et organisation}

La question demande de lire un fichier DIMACS dans une structure de données
Python, d'implémenter les algorithmes 1 et 2 de l'énoncé, de les valider puis
de commenter leurs performances. Le code se répartit en deux fichiers :
\begin{itemize}
  \item \texttt{dimacs.py} : lecture d'un fichier DIMACS et vérification
    qu'une valuation satisfait une formule ;
  \item \texttt{dpll.py} : algorithme 1 (\texttt{dpll\_sat}), algorithme 2
    (\texttt{dpll}) et fonctions auxiliaires.
\end{itemize}
Les tests se trouvent dans le dossier \texttt{tests} et s'exécutent avec :
\begin{lstlisting}
uv run python -m unittest discover -s tests -t .
\end{lstlisting}
Comme pour la question 1, seule la bibliothèque standard de Python 3.12 est
utilisée : c'est la version imposée pour la compétition de la question 6.

\section{Représentation des formules et lecture DIMACS}

\subsection{Structure de données}

Une formule CNF est représentée par un couple \texttt{(nb\_vars, clauses)}.
La liste \texttt{clauses} contient une clause par élément ; chaque clause est
un \texttt{frozenset} d'entiers. Nous conservons la convention DIMACS :
l'entier $k$ représente $x_k$ et l'entier $-k$ représente $\neg x_k$.
Par exemple, l'énigme des trois lampes de l'énoncé devient :
\begin{lstlisting}
(3, [frozenset({1, 2, 3}), frozenset({-1, -2}),
     frozenset({-2, -3}), frozenset({2, 3})])
\end{lstlisting}

Ce choix est guidé par l'opération centrale de DPLL, la simplification
$F[l\leftarrow 1]$ :
\begin{itemize}
  \item \textbf{Pas de traduction.} Les littéraux du fichier sont utilisés
    tels quels ; la négation d'un littéral \texttt{l} est simplement
    \texttt{-l} et sa variable \texttt{abs(l)}.
  \item \textbf{Tests d'appartenance en temps constant.} Savoir si une clause
    contient $l$ ou $\neg l$ est un test \texttt{in} sur un ensemble.
  \item \textbf{Doublons éliminés.} Une clause \texttt{1 1 -2 0} devient
    $\{1,-2\}$, ce qui correspond à la sémantique d'une clause, vue comme
    un ensemble de littéraux.
  \item \textbf{Immuabilité.} Un \texttt{frozenset} ne peut pas être modifié.
    Chaque appel récursif travaille sur une nouvelle formule et le retour
    arrière consiste simplement à revenir à la formule de l'appel parent,
    qui n'a pas été touchée. C'est exactement le modèle des algorithmes de
    l'énoncé, où $F[p\leftarrow b]$ et $F[p\leftarrow\overline{b}]$ sont
    deux formules indépendantes.
\end{itemize}

\paragraph{Alternatives écartées.} Une liste de listes est plus compacte,
mais retirer un littéral d'une clause y coûte un parcours de la clause et les
doublons doivent être gérés à part. Une structure mutable, avec annulation des
modifications lors du retour arrière (pile de modifications, littéraux
surveillés), est bien plus rapide : elle évite de recopier la formule à chaque
appel. Elle est en revanche plus délicate à écrire et à justifier. Nous la
réservons à la version optimisée des questions 4 à 6 ; la version de la
question 2 sert de référence simple et sûre.

\subsection{Règles de lecture}

La fonction \texttt{parse\_dimacs} lit le texte ligne par ligne :
\begin{itemize}
  \item les lignes vides et les lignes commençant par \texttt{c} sont ignorées ;
  \item l'en-tête \texttt{p cnf V C} est obligatoire, unique, et doit
    précéder toute clause ;
  \item une clause est une suite d'entiers terminée par \texttt{0}. Elle peut
    s'étendre sur plusieurs lignes, et plusieurs clauses peuvent partager
    une ligne : nous lisons un flux d'entiers plutôt qu'une clause par ligne ;
  \item une ligne \texttt{\%} arrête la lecture. Les fichiers de la
    bibliothèque SATLIB se terminent par \texttt{\%} puis \texttt{0}, qui
    serait sinon lu comme une clause vide et rendrait la formule
    insatisfiable à tort.
\end{itemize}
La fonction \texttt{read\_dimacs} applique \texttt{parse\_dimacs} au contenu
d'un fichier. Séparer les deux permet de tester la lecture sur des chaînes
de caractères, sans créer de fichiers.

\subsection{Erreurs et tolérances}

Comme pour \texttt{read\_meowdoku}, une entrée mal formée provoque une
\texttt{ValueError} avec le numéro de la ligne fautive : en-tête absent,
répété ou mal formé, littéral non entier, ou variable hors de
$\{1,\ldots,V\}$. Ce dernier contrôle est nécessaire, car le modèle affiché
à la question 6 porte sur les variables $1$ à $V$.

Deux écarts au format sont en revanche tolérés, car ils ne rendent pas la
formule ambiguë :
\begin{itemize}
  \item le nombre de clauses $C$ de l'en-tête n'est pas vérifié ; il n'est
    pas nécessaire à la résolution ;
  \item une dernière clause sans \texttt{0} final est conservée.
\end{itemize}
Ce choix a été fait avant de disposer des fichiers des enseignants : mieux
vaut qu'un fichier légèrement négligé soit lu correctement que rejeté, surtout
lors de la compétition où une instance non lue est une instance perdue.

Deux cas particuliers sont conservés tels quels. Une clause vide
(\texttt{0} seul) est gardée : DPLL répond alors immédiatement UNSAT, ce qui
est la bonne réponse. Une clause tautologique comme $x_1\lor\neg x_1$ est
également gardée : elle ne gêne pas DPLL et disparaît dès que $x_1$ est
affectée. La supprimer dès la lecture est une optimisation possible pour la
suite.

\subsection{Vérification d'un modèle}

La fonction \texttt{check\_model(clauses, model)} renvoie vrai si chaque clause
contient un littéral rendu vrai par la valuation \texttt{model}, un
dictionnaire \{variable : booléen\}. Un littéral dont la variable n'est pas
affectée n'est \emph{pas} considéré comme vrai. Cette définition stricte est
celle qu'appliquera le correcteur automatique de la compétition ; elle oblige
le solveur à fournir une valuation complète (voir
section~\ref{sec:valuation}). La fonction est indépendante du solveur, ce qui
en fait un juge impartial dans les tests.

\section{Implémentation de DPLL}

\subsection{Simplification}

La fonction \texttt{simplify(clauses, l)} calcule $F[l\leftarrow 1]$ en un
seul parcours, conformément à la définition de l'énoncé : une clause contenant
$l$ est vraie et disparaît ; une clause contenant $\neg l$ perd ce littéral et
peut devenir vide ; les autres clauses sont reprises sans copie. Affecter la
variable $p$ à la valeur $b$ revient à affecter le littéral $p$ (si $b=1$) ou
$-p$ (si $b=0$) à vrai : une seule fonction suffit donc pour les règles
unitaire, pure et de branchement.

\subsection{Correspondance avec les algorithmes de l'énoncé}

Les deux algorithmes appliquent les règles dans l'ordre de l'énoncé :
\begin{enumerate}
  \item clause vide : UNSAT ;
  \item formule sans clause : SAT ;
  \item clause unitaire $\{l\}$ : on continue avec $F[l\leftarrow 1]$ ;
  \item littéral pur $l$ : on continue avec $F[l\leftarrow 1]$ ;
  \item sinon, choix aléatoire $(p,b)$, essai de $F[p\leftarrow b]$, puis
    en cas d'échec de $F[p\leftarrow\overline{b}]$.
\end{enumerate}
Les étapes 1 à 4 sont communes aux deux versions et regroupées dans
\texttt{\_propagate}. \texttt{dpll\_sat} (algorithme 1) renvoie un booléen ;
\texttt{dpll} (algorithme 2) renvoie \texttt{(SAT, v)} ou
\texttt{(UNSAT, None)}. Nous utilisons \texttt{None} plutôt que le $0$ de
l'énoncé : c'est la convention Python pour \og pas de valeur \fg{} et elle ne
peut pas être confondue avec une valuation.

Un littéral pur est cherché en rassemblant tous les littéraux de la formule
dans un ensemble, puis en prenant le premier $l$ (par numéro de variable
croissant) tel que $\neg l$ est absent. Le tri rend le choix indépendant de
l'ordre interne de l'ensemble, donc reproductible.

\subsection{Propagation par une boucle plutôt que par récursion}

Dans l'énoncé, les règles unitaire et pure sont des appels récursifs : chaque
clause unitaire ajoute un niveau de récursion. Or Python limite par défaut la
profondeur de récursion à 1000 appels. Une formule qui propage plusieurs
centaines de clauses unitaires, ce qui est courant (une grille Meowdoku de
taille $n$ a $n^2$ variables), ferait échouer le programme.

Nous remplaçons donc ces deux appels par une boucle. La transformation ne
change pas le résultat : dans les lignes 5 à 16 de l'algorithme 2, l'appel
récursif est le dernier calcul effectué. Son statut est renvoyé tel quel et la
valuation est seulement complétée par $\{l\}$. Il suffit donc de mémoriser les
littéraux fixés dans une liste (\texttt{forced}) et de les ajouter à la
valuation une fois le statut connu.

Seul le branchement reste récursif, car il nécessite un vrai retour arrière.
La profondeur de récursion est alors au plus le nombre de décisions d'une
branche, donc au plus le nombre de variables : chaque décision fait
disparaître sa variable. Par sécurité, la limite de récursion est relevée au
nombre de variables plus 1000 si nécessaire. Depuis Python 3.11, les appels
entre fonctions Python ne consomment plus la pile du langage C, ce qui rend
cette augmentation sans risque. Un test vérifie une formule demandant 3000
décisions imbriquées.

\subsection{Choix aléatoire et reproductibilité}

La fonction \texttt{random\_choice} correspond à \texttt{RandomChoice} : elle
tire uniformément une variable parmi celles qui apparaissent encore dans la
formule, puis une valeur de vérité. Le couple $(p,b)$ est renvoyé sous la forme
d'un littéral ($p$ ou $-p$) ; la branche \og sinon \fg{} utilise son opposé.
Isoler ce choix dans une fonction prépare la question 4, où il sera remplacé
par \texttt{Heuristics(F, H)}.

Le générateur aléatoire est passé en paramètre (\texttt{rng}, un objet
\texttt{random.Random}). Avec une graine fixée, une exécution est donc
reproductible, ce qui est indispensable pour déboguer et pour comparer les
mesures de performance. Les variables sont triées avant le tirage pour que le
résultat ne dépende que de la graine.

\subsection{Valuation partielle et complétion}
\label{sec:valuation}

La valuation renvoyée par l'algorithme 2 est partielle. Une variable peut
disparaître de la formule sans être affectée : c'est le cas lorsque toutes les
clauses qui la contiennent ont été satisfaites par d'autres littéraux. Par
exemple, pour la formule $(x_1\lor x_2)$, le littéral $x_1$ est pur : la règle
fixe $x_1$ à vrai, la clause disparaît et la formule est vide. La valuation
renvoyée est $\{x_1\leftarrow 1\}$ et $x_2$ n'y figure pas.

Toute valeur convient alors pour ces variables : au moment où elles
disparaissent, chaque clause qui les contient est déjà satisfaite par un
littéral affecté. La fonction \texttt{complete} leur donne la valeur faux
et produit une valuation sur toutes les variables $1,\ldots,V$, ce qu'exigent
\texttt{check\_model} et le format de sortie de la question 6.

\subsection{Compteurs}

Une structure \texttt{Stats} compte les décisions, les retours arrière, les
propagations unitaires et les applications de la règle du littéral pur. Ces
compteurs correspondent aux critères de la question 5. Leur coût est
négligeable (une incrémentation par étape) et ils permettent dès maintenant
d'expliquer les temps observés.

\subsection{Complexité}

Notons $|F|$ la taille de la formule, c'est-à-dire le nombre total de
littéraux. Chaque n\oe{}ud de l'arbre de recherche coûte $O(|F|)$ pour
\texttt{simplify}, la recherche de clause vide et de clause unitaire, et
$O(|F| + v\log v)$ pour la recherche de littéral pur et le choix aléatoire,
avec $v$ le nombre de variables restantes. Une propagation de $u$ littéraux
coûte donc $O(u\,|F|)$. Le nombre de n\oe{}uds est au pire exponentiel en le
nombre de variables : DPLL reste exponentiel dans le pire cas, ce qui est
attendu pour un problème NP-complet.

\section{Validation}

La validation repose sur plusieurs sources indépendantes : un oracle par
force brute, des grilles Meowdoku, des cas limites, des benchmarks SATLIB et
les instances publiques fournies par les enseignants.

\paragraph{Oracle par force brute.} 1500 formules aléatoires de 1 à 10
variables, avec jusqu'à $5n$ clauses de 1 à 4 littéraux, sont résolues par les
deux algorithmes et par une énumération des $2^n$ valuations. Le jeu contient
plusieurs centaines d'instances de chaque statut, ce que le test vérifie.
Pour chaque formule, les deux algorithmes doivent donner le statut de
l'oracle, et toute valuation renvoyée, une fois complétée, doit satisfaire
la formule d'après \texttt{check\_model}.

Pour nous assurer que ce test détecte réellement les erreurs, nous avons
volontairement introduit un défaut dans \texttt{simplify} (oubli de retirer
$\neg l$ des clauses) : 979 des 1500 cas échouent alors.

\paragraph{Instances Meowdoku.} Les grilles générées à la question 1 pour
$n=4$ à $9$ sont satisfiables par construction ; le solveur doit répondre SAT
avec un modèle valide. Trois grilles construites à la main doivent être
insatisfiables : une grille $2\times2$ (deux chats forcément voisins), une
grille $3\times3$ (aucune solution n'existe pour $n=3$) et une grille
$4\times4$ où deux régions d'une seule case sont adjacentes.

\paragraph{Cas limites.} Formule vide (SAT), clause vide (UNSAT),
$x\land\neg x$ (UNSAT), tautologie (SAT), exemple des trois lampes,
chaîne de 3000 implications se terminant par une contradiction (uniquement
de la propagation unitaire) et 3000 décisions imbriquées.

\paragraph{Benchmarks SATLIB.} Avant de recevoir les fichiers des
enseignants, nous avons utilisé les instances 3-SAT aléatoires de la bibliothèque de référence
SATLIB\footnote{\url{https://www.cs.ubc.ca/~hoos/SATLIB/benchm.html}},
placées dans \texttt{data/satlib} par le script \texttt{download\_satlib.py}
(elles ne sont pas incluses dans le dépôt). Leur statut est connu par
construction :
les familles \texttt{uf} sont satisfiables, les familles \texttt{uuf}
insatisfiables. Ces fichiers se terminent par la ligne \texttt{\%} suivie
d'un \texttt{0} : ils ont confirmé l'utilité de la règle d'arrêt du lecteur.

Le script \texttt{validate.py} parcourt un dossier et, pour chaque fichier,
lance les deux algorithmes avec la même graine. Il vérifie que les deux
statuts concordent, que le statut attendu (lu dans un fichier
\texttt{expected.json} s'il existe, sinon déduit du nom : préfixe
\texttt{uf} ou \texttt{uuf}, ou mention \texttt{sat} / \texttt{unsat}) est
obtenu, et que le modèle complété satisfait la formule. Chaque algorithme
est interrompu après 60~s, la limite de la compétition :
\begin{lstlisting}
uv run python validate.py data/satlib/uf50-218
\end{lstlisting}

\begin{center}
\begin{tabular}{lrrrrr}
  Famille & Variables & Clauses & Instances & Échecs & Temps total (algos 1 et 2) \\
  \hline
  \texttt{uf20-91}   & 20 & 91  & 1000 & 0 & 1,0 s \\
  \texttt{uf50-218}  & 50 & 218 & 1000 & 0 & 52,9 s \\
  \texttt{uuf50-218} & 50 & 218 & 1000 & 0 & 147,7 s \\
  \texttt{uf75-325}  & 75 & 325 & 100  & 0 & 104,3 s \\
  \texttt{uuf75-325} & 75 & 325 & 100  & 0 & 264,3 s \\
\end{tabular}
\end{center}
Les 3200 instances sont correctement résolues par les deux algorithmes,
sans aucun dépassement du délai. Le temps indiqué est le total pour toute la
famille : pour uuf75-325, 264{,}3~s pour 100 instances, soit environ 2{,}6~s
par instance pour les deux algorithmes, bien en dessous de la limite de
60~s par instance.

\paragraph{Instances publiques des enseignants.} Le dossier
\texttt{public-instances} contient 41 fichiers et un fichier
\texttt{expected.json} donnant le statut attendu de chacun : 11 petits tests
d'interface et cas limites (formule vide, formule sans clause, clause vide,
tautologie, littéraux répétés, littéraux tous négatifs), et 30 instances de
benchmark issues de SATLIB (19 SAT, 11 UNSAT, de 42 à 1363 variables).
\begin{lstlisting}
uv run python validate.py public-instances -v --timeout 60
\end{lstlisting}

\begin{center}
\begin{tabular}{lrrrl}
  Instance & Variables & Clauses & Attendu & Résultat (graine 0, algos 1 et 2) \\
  \hline
  \texttt{public-11} & 729 & 21844 & SAT   & SAT en 29,8 s \\
  \texttt{public-18} & 459 & 7054  & SAT   & SAT en 9,9 s \\
  \texttt{public-19} & 729 & 22060 & SAT   & SAT en 10,3 s \\
  \texttt{public-03} & 63  & 168   & UNSAT & timeout (60 s) \\
  \texttt{public-12} & 522 & 11636 & SAT   & timeout (60 s) \\
  \texttt{public-27} & 181 & 3151  & SAT   & timeout (60 s) \\
  27 autres benchmarks & & & & tous corrects, moins de 2 s chacun \\
  11 tests d'interface & & & & tous corrects, instantanés \\
\end{tabular}
\end{center}

Les 38 instances résolues dans le délai donnent toutes le statut attendu,
avec des modèles valides et des algorithmes 1 et 2 concordants. Aucune
réponse n'est fausse. Les 11 petits fichiers sont également testés à travers
l'interface \texttt{sat.py}, dans les deux modes, exactement comme le fera
le correcteur (fichier \texttt{tests/test\_sat.py}). Les cas limites étaient
déjà tous couverts par nos propres tests, ce qui confirme les choix du
lecteur DIMACS (littéraux répétés, clause vide, formule sans clause).

Trois instances dépassent 60~s avec la graine 0. Pour savoir si la
difficulté vient de l'instance ou du hasard, nous les avons relancées avec
les graines 1, 2 et 3 (algorithme 2 seul) :
\begin{center}
\begin{tabular}{lrrr}
  Instance & Graine 1 & Graine 2 & Graine 3 \\
  \hline
  \texttt{public-12} (SAT)   & 1,1 s    & 2,5 s    & timeout \\
  \texttt{public-27} (SAT)   & 2,5 s    & 53,0 s   & 6,8 s \\
  \texttt{public-03} (UNSAT) & timeout  & timeout  & timeout \\
\end{tabular}
\end{center}
Les deux instances satisfiables sont résolues en quelques secondes avec
une graine favorable : comme sur Meowdoku, c'est la queue lourde du choix
aléatoire. L'instance \texttt{public-03} est d'une autre nature : elle est
insatisfiable et petite (63 variables, 168 clauses de 3 littéraux), mais
aucune graine ne la résout. En 60~s, DPLL y prend environ 1,8 million de
décisions, toutes suivies d'un retour arrière, sans terminer. Il s'agit
vraisemblablement d'une des deux instances annoncées comme volontairement
difficiles. Une heuristique ne suffira probablement pas : ce type
d'instance motive des techniques qui dépassent DPLL (apprentissage de
clauses). Nous y reviendrons à la question 6.

\paragraph{Interface en ligne de commande.} Le fichier \texttt{sat.py}
respecte dès maintenant l'interface imposée à la question 6 :
\begin{lstlisting}
python sat.py instance.cnf --mode dpll --output-model
\end{lstlisting}
Le mode \texttt{dpll} utilise l'algorithme 2 ; le mode \texttt{optimized}
l'utilise aussi en attendant la question 6. La sortie standard ne contient
que la ligne \texttt{s SATISFIABLE} ou \texttt{s UNSATISFIABLE} et, sur
demande, la ligne \texttt{v} du modèle complet, terminée par \texttt{0}. Un
fichier illisible produit un message sur la sortie d'erreur et un code de
retour non nul, sans rien écrire sur la sortie standard : une réponse
incorrecte serait pénalisée en compétition. Des tests lancent le programme
dans un processus séparé et vérifient exactement ce format.

\section{Performances}

\subsection{Protocole}

Le script \texttt{bench\_q2.py} mesure le temps de résolution
(\texttt{time.perf\_counter}) et relève les compteurs de \texttt{Stats} :
décisions, retours arrière, propagations unitaires et littéraux purs. Chaque
résolution est enregistrée dans \texttt{results/q2\_<expérience>.csv} ;
les tableaux ci-dessous donnent des \emph{médianes}, peu sensibles aux
quelques exécutions extrêmes que le choix aléatoire produit. Une résolution
est interrompue au bout de 30~s et comptée comme \og timeout \fg{}. Les
mesures ont été faites sur un Mac (Apple M1) avec Python 3.12 ; seuls les
ordres de grandeur et les évolutions comptent, pas les valeurs absolues.
\begin{lstlisting}
uv run python bench_q2.py --timeout 30
\end{lstlisting}

Nous utilisons cinq familles d'instances :
\begin{itemize}
  \item des instances SATLIB, résolues par les deux algorithmes avec trois
    graines ;
  \item du 3-SAT aléatoire à $n=40$ variables, en faisant varier le ratio
    $m/n$ (nombre de clauses sur nombre de variables) ;
  \item du 3-SAT aléatoire au ratio $4{,}26$, en faisant varier $n$ ;
  \item des grilles Meowdoku générées par la question 1 ;
  \item une même instance résolue avec de nombreuses graines.
\end{itemize}
Le 3-SAT aléatoire suit le même modèle que SATLIB : chaque clause porte sur
trois variables distinctes tirées au hasard, avec des signes aléatoires.

\subsection{Algorithme 1 contre algorithme 2}

\begin{center}
\begin{tabular}{lrrrrrr}
  Famille & Runs & Algo & Temps méd. & Temps max & Déc. méd. & Unit. méd. \\
  \hline
  \texttt{uf50-218}  & 150 & 1 & 0,020 s & 0,098 s & 133 & 1594 \\
                     & 150 & 2 & 0,020 s & 0,099 s & 133 & 1594 \\
  \texttt{uuf50-218} & 150 & 1 & 0,072 s & 0,131 s & 461 & 5555 \\
                     & 150 & 2 & 0,072 s & 0,131 s & 461 & 5555 \\
  \texttt{uf75-325}  & 90  & 1 & 0,352 s & 2,12 s  & 1336 & 21082 \\
                     & 90  & 2 & 0,349 s & 2,13 s  & 1336 & 21082 \\
  \texttt{uuf75-325} & 90  & 1 & 1,29 s  & 2,25 s  & 4696 & 73920 \\
                     & 90  & 2 & 1,32 s  & 2,29 s  & 4696 & 73920 \\
\end{tabular}
\end{center}

Avec la même graine, les deux algorithmes font exactement les mêmes choix et
parcourent le même arbre : tous les compteurs sont identiques. Les temps
diffèrent de moins de 3~\% : construire la valuation (une écriture dans un
dictionnaire par variable affectée, sur la seule branche gagnante) est
négligeable devant les simplifications. L'algorithme 2 ne coûte donc rien
de plus et c'est lui que nous utilisons partout ailleurs.

\subsection{Effet du ratio clauses / variables}

\begin{center}
\begin{tabular}{rrrrrrr}
  $m/n$ & SAT & UNSAT & Temps méd. & Décisions méd. & Unitaires méd. & Purs méd. \\
  \hline
  2    & 40 & 0  & 0,4 ms  & 7   & 4    & 26 \\
  3    & 40 & 0  & 0,5 ms  & 11  & 30   & 11 \\
  3,5  & 40 & 0  & 0,9 ms  & 16  & 90   & 11 \\
  4    & 32 & 8  & 4,8 ms  & 48  & 458  & 15 \\
  4,26 & 26 & 14 & 9,2 ms  & 86  & 853  & 15 \\
  4,5  & 22 & 18 & 12,3 ms & 105 & 1080 & 11 \\
  5    & 3  & 37 & 15,4 ms & 106 & 1104 & 3 \\
  6    & 0  & 40 & 10,6 ms & 60  & 604  & 0 \\
  8    & 0  & 40 & 7,3 ms  & 26  & 255  & 0 \\
\end{tabular}
\end{center}
\noindent($n=40$, 40 formules par ratio.)

On retrouve le phénomène de \emph{transition de phase} du 3-SAT aléatoire.
Avec peu de clauses, presque toute valuation convient : DPLL trouve une
solution quasiment sans retour arrière, et la règle du littéral pur est très
utile (26 applications pour 7 décisions à $m/n=2$). Avec beaucoup de clauses,
la formule est très contrainte : la propagation unitaire produit vite une
clause vide et l'insatisfiabilité est prouvée par un petit arbre. La
difficulté est maximale entre les deux, là où la proportion d'instances
satisfiables passe de 1 à 0, autour de $m/n\approx 4{,}26$ (une moitié
environ de SAT). Pour $n=40$, le pic de coût se situe un peu au-delà, vers
4,5--5 : à petite taille, la transition est étalée et les instances UNSAT
proches du seuil, plus coûteuses, tirent la médiane vers la droite.

La règle du littéral pur devient inutile quand $m/n$ augmente : chaque
variable apparaît dans environ $3m/n$ clauses, et il est de moins en moins
probable qu'elle n'apparaisse que sous une seule polarité.

\subsection{Croissance avec le nombre de variables}

\begin{center}
\begin{tabular}{rrrrrrr}
  $n$ & Temps méd. & Temps max & Déc. méd. & Déc. SAT & Déc. UNSAT & ms / déc. \\
  \hline
  20 & 0,7 ms  & 2,0 ms  & 15   & 8    & 21   & 0,047 \\
  30 & 4,1 ms  & 8,2 ms  & 51   & 39   & 83   & 0,078 \\
  40 & 11 ms   & 50 ms   & 96   & 24   & 211  & 0,111 \\
  50 & 37 ms   & 84 ms   & 247  & 137  & 438  & 0,150 \\
  60 & 160 ms  & 563 ms  & 856  & 327  & 1570 & 0,185 \\
  70 & 635 ms  & 1447 ms & 2632 & 1086 & 3780 & 0,235 \\
\end{tabular}
\end{center}
\noindent(Ratio 4,26, 20 formules par taille, environ moitié SAT et moitié
UNSAT.)

\paragraph{Croissance exponentielle.} De $n=20$ à $n=70$, le nombre médian de
décisions est multiplié par environ 175, soit un facteur
$175^{1/50}\approx 1{,}11$ par variable supplémentaire : le nombre de
décisions double environ toutes les 7 variables. C'est la croissance exponentielle attendue
de DPLL sur des instances difficiles. À ce rythme, une instance aléatoire
de 150 variables au seuil demanderait plus d'une heure.

\paragraph{Coût d'une décision.} Le temps moyen par décision augmente aussi,
de 0,05 à 0,24~ms, à peu près proportionnellement à $n$. C'est la
conséquence directe de notre structure de données : chaque simplification
recopie la formule, de taille proportionnelle à $m = 4{,}26\,n$, et une
décision entraîne en moyenne une quinzaine de propagations unitaires, donc
une quinzaine de copies. La complexité observée est bien
$(\text{nombre de n\oe{}uds}) \times O(|F|)$.

\paragraph{SAT contre UNSAT.} À taille égale, les instances UNSAT demandent
2 à 9 fois plus de décisions que les SAT. Pour prouver l'insatisfiabilité,
DPLL doit réfuter \emph{toutes} les branches, alors qu'il s'arrête dès la
première solution d'une formule satisfiable. Sur les instances UNSAT, le
nombre de retours arrière est d'ailleurs exactement égal au nombre de
décisions : la première branche de chaque décision échoue forcément, sinon
la formule serait satisfiable. C'est aussi pourquoi les temps des instances
UNSAT sont plus réguliers (temps maximum proche de la médiane) : l'arbre
complet dépend moins de la chance que la recherche d'une seule solution.

\subsection{Instances Meowdoku}

\begin{center}
\begin{tabular}{rrrrrrrr}
  $n$ & Var. & Résolues & Timeouts & Temps méd. & Temps max & Déc. méd. & Purs \\
  \hline
  4  & 16  & 10 & 0 & 0,2 ms  & 0,2 ms & 5    & 0 \\
  6  & 36  & 10 & 0 & 0,7 ms  & 2,8 ms & 6    & 0 \\
  8  & 64  & 10 & 0 & 1,8 ms  & 7,4 ms & 10   & 0 \\
  10 & 100 & 10 & 0 & 4,5 ms  & 86 ms  & 13   & 0 \\
  12 & 144 & 10 & 0 & 17 ms   & 0,75 s & 48   & 0 \\
  14 & 196 & 9  & 1 & 34 ms   & 26 s   & 177  & 0 \\
  16 & 256 & 6  & 4 & 1,3 s   & 7,8 s  & 6178 & 0 \\
\end{tabular}
\end{center}
\noindent(10 grilles par taille, graine 0 ; médiane et maximum calculés sur
les grilles résolues avant 30~s.)

Jusqu'à $n=10$, les grilles sont résolues en quelques millisecondes : la
structure du problème est très favorable à la propagation unitaire. Dès
qu'un chat est placé, toutes les cases de sa ligne, de sa colonne, de sa
région et ses voisines diagonales sont mises à faux par propagation, et
une région réduite à une seule case libre force immédiatement son chat. Une
dizaine de décisions suffit pour 100 variables, bien moins que pour du 3-SAT
aléatoire de même taille.

La règle du littéral pur ne s'applique \emph{jamais}, sur aucune grille.
Cela s'explique par l'encodage. Une variable encore présente appartient à une
région dont la clause positive n'est pas satisfaite (sinon la propagation
l'aurait déjà mise à faux, par les clauses \og au plus un chat par région
\fg{}) : elle apparaît donc positivement. Elle apparaît aussi négativement
dans les clauses binaires de sa ligne, de sa colonne ou de ses diagonales.
Elle a donc les deux polarités, et n'est pas pure. Sur Meowdoku, la
recherche de littéraux purs est un coût sans aucun bénéfice : c'est une
piste directe d'optimisation pour la question 6.

À partir de $n=12$, le comportement devient très irrégulier : la plupart
des grilles restent faciles, mais certaines demandent des centaines de
milliers de décisions. À $n=16$, 4 grilles sur 10 dépassent 30~s. Le choix
aléatoire n'exploite pas la structure du problème : il peut par exemple
choisir de mettre à faux une case quelconque, ce qui ne propage presque
rien, au lieu de placer un chat dans une région presque pleine.

\subsection{Sensibilité au hasard}

\begin{center}
\begin{tabular}{lrrrrr}
  Instance & Graines & Temps min & Temps méd. & Temps max & Timeouts \\
  \hline
  \texttt{uf50-01} (SAT)    & 100 & 0,001 s & 0,035 s & 0,131 s & 0 \\
  \texttt{uuf50-01} (UNSAT) & 100 & 0,040 s & 0,065 s & 0,100 s & 0 \\
  Meowdoku $n=14$ (grille 2) & 20 & 0,020 s & 0,061 s & $>30$ s & 4 \\
\end{tabular}
\end{center}

Une même formule peut être résolue 100 fois plus vite selon la graine
lorsqu'elle est satisfiable : une suite de bons choix mène directement à une
solution. Pour l'instance UNSAT, l'écart est bien plus faible (facteur 2,5),
car l'arbre doit de toute façon être exploré en entier ; seul son nombre de
n\oe{}uds dépend de l'ordre des choix. Le cas Meowdoku est le plus frappant :
la grille qui avait dépassé le délai dans l'expérience précédente est résolue
en 0,06~s pour la plupart des graines, mais 4 graines sur 20 dépassent
encore 30~s. La distribution du temps a une \og queue lourde \fg{} : un seul
mauvais choix au début de la recherche peut coûter une exploration énorme.

\subsection{Bilan}

\begin{itemize}
  \item Les deux algorithmes ont les mêmes performances ; l'algorithme 2
    fournit le modèle sans surcoût mesurable.
  \item Le coût suit $(\text{nombre de décisions}) \times O(|F|)$ : le
    nombre de décisions croît exponentiellement sur les instances
    difficiles, et chaque décision coûte une copie de la formule pour chaque
    simplification.
  \item Les instances UNSAT sont plus coûteuses et plus régulières que les
    instances SAT.
  \item Sur Meowdoku, la propagation unitaire fait l'essentiel du travail, la
    règle du littéral pur est inutile, et le choix aléatoire provoque de
    rares mais énormes explosions du temps de calcul à partir de $n\approx 14$.
\end{itemize}
Ces observations orientent la suite : des heuristiques de branchement
(question 4) pour choisir des variables qui déclenchent beaucoup de
propagation, puis une structure de données sans recopie et la suppression de
la recherche de littéraux purs pour la version optimisée (question 6).

\end{document}
